%% ============================================================
%%  CHAPTER 4 — METHODOLOGY OVERVIEW
%% ============================================================
\chapter{Methodology Overview}
\label{ch:methodology}

This chapter provides a high-level, end-to-end walkthrough of the
complete methodology before the subsequent chapters examine each
stage in detail. The aim is to give the reader a clear mental map
of how a raw \acrshort{eeg} recording is progressively transformed
into a physical prosthetic hand action, and how the individual
processing stages connect to one another. Each stage is introduced
conceptually and illustrated with a diagram; the full mathematical,
algorithmic, and implementation details are deferred to the
dedicated chapters that follow.

Section~\ref{sec:mo_master} presents the master architecture of the
complete system. Section~\ref{sec:mo_dataflow} traces the
transformation of the data through the pipeline, showing how its
shape and meaning change at each stage.
Section~\ref{sec:mo_collection} introduces the data collection
stage. Section~\ref{sec:mo_preprocessing} introduces the signal
processing stage. Section~\ref{sec:mo_features} introduces the
feature extraction stage. Section~\ref{sec:mo_classification}
introduces the classification stage. Section~\ref{sec:mo_inference}
introduces the real-time inference stage.
Section~\ref{sec:mo_summary} summarises the chapter and maps each
stage to its detailed chapter.

%% ----------------------------------------------------------
\section{Master System Architecture}
\label{sec:mo_master}
%% ----------------------------------------------------------

At the highest level, the system is a pipeline that converts brain
and muscle electrical activity --- captured at the scalp --- into
discrete control commands for a prosthetic hand. The pipeline
operates in two distinct modes that share the same core processing
stages: an offline \emph{training mode} that learns a classifier
from pre-recorded data, and an online \emph{inference mode} that
applies the learned classifier to a live signal stream in real
time.

Figure~\ref{fig:mo_master} shows the master architecture. The
shared stages --- preprocessing, activity gating, and feature
extraction --- are deliberately identical in both modes, which is
the key to ensuring that the model encounters the same kind of data
during deployment as it did during training.

\begin{figure}[htbp]
\centering
\begin{tikzpicture}[
    node distance=0.55cm and 1.1cm,
    box/.style={rectangle, rounded corners=2pt, draw=black!70,
                very thick, minimum width=3.0cm, minimum height=0.9cm,
                align=center, font=\footnotesize},
    shared/.style={box, fill=blue!12},
    train/.style={box, fill=orange!15},
    infer/.style={box, fill=green!12},
    hw/.style={box, fill=red!12},
    arow/.style={-{Stealth[length=2mm]}, thick},
]

% Acquisition (shared top)
\node[shared] (acq) {Stage 1: EEG Acquisition\\\scriptsize 19 ch, 200\,Hz, EDF};

% Two modes branching
\node[train, below left=0.9cm and -1.0cm of acq] (edf) {Training Mode\\\scriptsize Recorded EDF files};
\node[infer, below right=0.9cm and -1.0cm of acq] (stream) {Inference Mode\\\scriptsize Live EEG stream};

% Shared preprocessing
\node[shared, below=1.5cm of acq] (pre) {Stage 2: Preprocessing\\\scriptsize Notch, bandpass, CAR};
\node[shared, below=0.55cm of pre] (gate) {Stage 3: Activity Gate\\\scriptsize RMS threshold};
\node[shared, below=0.55cm of gate] (feat) {Stage 4: Feature Extraction\\\scriptsize channel-aware features};

% Split: train vs classify
\node[train, below left=0.9cm and 0.1cm of feat] (train) {Train Hybrid Model\\\scriptsize features + DTW + Riemann};
\node[infer, below right=0.9cm and 0.1cm of feat] (clf) {Stage 5: Classify\\\scriptsize hybrid router};

% Model store
\node[box, fill=gray!15, below=0.7cm of train] (model) {Model store\\\scriptsize trained hybrid model};

% Inference engine + hardware
\node[infer, below=0.9cm of clf] (fsm) {Stage 6: Inference Engine\\\scriptsize Confidence + FSM};
\node[hw, below=0.55cm of fsm] (hand) {Pi5 $\to$ ESP32 $\to$ hand\\\scriptsize label $\to$ 6 servos};

% Arrows: acquisition to modes
\draw[arow] (acq) -- (edf);
\draw[arow] (acq) -- (stream);
% modes into preprocessing
\draw[arow] (edf) -- (pre);
\draw[arow] (stream) -- (pre);
% shared chain
\draw[arow] (pre) -- (gate);
\draw[arow] (gate) -- (feat);
% feat to train and classify
\draw[arow] (feat) -- (train);
\draw[arow] (feat) -- (clf);
% train to model to classify
\draw[arow] (train) -- (model);
\draw[arow] (model.east) -- ++(0.6,0) |- (clf.west);
% classify to fsm to hand
\draw[arow] (clf) -- (fsm);
\draw[arow] (fsm) -- (hand);

\end{tikzpicture}
\caption{Master system architecture. Blue stages are shared between
training and inference, ensuring train--inference consistency.
Orange stages are training-only; green stages are inference-only;
the red stage is the physical hardware. The hybrid model learned in
training mode is loaded by the inference engine, which emits an
abstract command label to the ESP32 hand controller for live
operation.}
\label{fig:mo_master}
\end{figure}

The single most important architectural decision visible in
Figure~\ref{fig:mo_master} is that Stages~2--4 (preprocessing,
activity gating, and feature extraction) are reused verbatim in
both modes. A classifier trained on features computed one way will
fail if, at inference time, the features are computed even slightly
differently. By sharing the exact same code path, the system
guarantees that an epoch seen during live inference is processed
identically to the epochs the model was trained on.

%% ----------------------------------------------------------
\section{Data Transformation Through the Pipeline}
\label{sec:mo_dataflow}
%% ----------------------------------------------------------

As data flows through the pipeline, its representation changes at
each stage --- from a long continuous multichannel recording, to
short windowed epochs, to a compact numerical feature vector, to a
single class label, and finally to a one-byte command character
sent to the prosthetic hand. Figure~\ref{fig:mo_dataflow} traces
this transformation, making explicit the shape and meaning of the
data at each step.

\begin{figure}[htbp]
\centering
\begin{tikzpicture}[
    node distance=0.9cm,
    stage/.style={rectangle, rounded corners=2pt, draw=black!60,
                  thick, minimum width=3.3cm, minimum height=0.8cm,
                  align=center, font=\scriptsize, fill=blue!8},
    data/.style={rectangle, draw=black!30, dashed, minimum width=4.0cm,
                 minimum height=0.7cm, align=center,
                 font=\scriptsize\ttfamily, fill=gray!8},
    arow/.style={-{Stealth[length=2mm]}, thick},
]

\node[stage] (s1) {Raw recording};
\node[data, right=1.5cm of s1] (d1) {19 $\times$ N samples};

\node[stage, below=0.9cm of s1] (s2) {Preprocess + epoch};
\node[data, right=1.5cm of s2] (d2) {19 $\times$ 400 per epoch};

\node[stage, below=0.9cm of s2] (s3) {Activity gate};
\node[data, right=1.5cm of s3] (d3) {active epochs only};

\node[stage, below=0.9cm of s3] (s4) {Feature extraction};
\node[data, right=1.5cm of s4] (d4) {feature vector};

\node[stage, below=0.9cm of s4] (s5) {Hybrid classify};
\node[data, right=1.5cm of s5] (d5) {1 of 5 classes};

\node[stage, below=0.9cm of s5] (s6) {FSM + UART};
\node[data, right=1.5cm of s6] (d6) {1 command label};

\foreach \a/\b in {s1/s2, s2/s3, s3/s4, s4/s5, s5/s6}
    \draw[arow] (\a) -- (\b);

\foreach \s/\d in {s1/d1, s2/d2, s3/d3, s4/d4, s5/d5, s6/d6}
    \draw[arow, gray] (\s) -- (\d);

\end{tikzpicture}
\caption{Data transformation through the pipeline. The left column
shows the processing stage; the right column (dashed) shows the
shape and meaning of the data after that stage. A continuous
multichannel recording is progressively reduced to a single command
label sent to the hand controller.}
\label{fig:mo_dataflow}
\end{figure}

This progressive reduction --- from raw signal to command --- is the
essence of the system. Each stage discards information that is
irrelevant to the control task while preserving and concentrating
the information that distinguishes one gesture from another. The
following sections introduce each stage in turn.

%% ----------------------------------------------------------
\section{Stage 1--2: Data Collection and Signal Acquisition}
\label{sec:mo_collection}
%% ----------------------------------------------------------

The pipeline begins with the deliberate production of voluntary
\acrshort{eeg} artifacts by the subject. Rather than relying on
subtle motor imagery, the subject performs five large, distinct
facial and jaw gestures --- single and double eye blinks; a
bilateral jaw clinch; and left and right lateral jaw bruxism ---
each of which produces a characteristic, high-amplitude signal at
specific electrode sites. A 19-channel \acrshort{eeg} device records
these signals at \SI{200}{\hertz} and stores them in
\acrshort{edf} files, one continuous recording per gesture class.

\begin{figure}[htbp]
\centering
\begin{tikzpicture}[
    node distance=0.8cm,
    src/.style={rectangle, rounded corners=2pt, draw=black!60, thick,
                minimum width=2.2cm, minimum height=0.8cm,
                align=center, font=\scriptsize},
    blink/.style={src, fill=blue!12},
    jaw/.style={src, fill=orange!15},
    sink/.style={rectangle, rounded corners=2pt, draw=black!70,
                 very thick, minimum width=2.4cm, minimum height=1.0cm,
                 align=center, font=\scriptsize, fill=gray!12},
    arow/.style={-{Stealth[length=2mm]}, thick},
]

\node[blink] (b1) {Single blink\\\scriptsize Fp1, Fp2};
\node[blink, below=0.4cm of b1] (b2) {Double blink\\\scriptsize Fp1, Fp2};
\node[jaw, below=0.4cm of b2] (j1) {Clinch\\\scriptsize T3 $\approx$ T4};
\node[jaw, below=0.4cm of j1] (j2) {Left bruxism\\\scriptsize T3 $>$ T4};
\node[jaw, below=0.4cm of j2] (j3) {Right bruxism\\\scriptsize T4 $>$ T3};

\node[sink, right=2.8cm of b2, yshift=-1.0cm] (edf) {19-channel\\EDF recording\\\scriptsize 200\,Hz};

\foreach \n in {b1,b2,j1,j2,j3}
    \draw[arow] (\n.east) -- (edf.west);

\end{tikzpicture}
\caption{Stage 1--2: five voluntary artifact gestures and the
electrode sites at which each is most prominent. Blink gestures
(blue) appear at the frontal electrodes Fp1/Fp2; jaw gestures
(orange) appear at the temporal electrodes T3/T4, with the
left/right amplitude balance distinguishing the three jaw classes.
All gestures are captured by the same 19-channel EDF recording.}
\label{fig:mo_collection}
\end{figure}

A central challenge introduced at this stage is that each recording
contains not only the active gesture events but also long rest
periods between them. Because the entire recording carries a single
class label, these rest windows become mislabelled training data ---
a problem resolved by the activity gate in Stage~3. The complete
data collection protocol, gesture-by-gesture physiology, recording
environment, and dataset statistics are described in
Chapter~6.

%% ----------------------------------------------------------
\section{Stage 2: Signal Processing}
\label{sec:mo_preprocessing}
%% ----------------------------------------------------------

The raw \acrshort{edf} recordings are not directly usable for
classification: they contain power-line interference, low-frequency
drift, broadband noise, and spatial common-mode signals. The signal
processing stage cleans each recording through a fixed sequence of
operations and then slices the continuous signal into short,
fixed-length windows (epochs) suitable for feature extraction.

\begin{figure}[htbp]
\centering
\begin{tikzpicture}[
    node distance=0.45cm,
    step/.style={rectangle, rounded corners=2pt, draw=black!60, thick,
                 minimum width=6.4cm, minimum height=0.7cm,
                 align=center, font=\scriptsize, fill=blue!8},
    arow/.style={-{Stealth[length=2mm]}, thick},
]

\node[step] (n)  {Notch filter \,---\, remove 50\,Hz power-line};
\node[step, below=of n]  (bp) {Bandpass 1--45\,Hz \,---\, keep blink + jaw EMG};
\node[step, below=of bp] (car){Common Average Reference \,---\, spatial clean};
\node[step, below=of car](ep) {Epoch \,---\, 2.0\,s windows, 75\% overlap};
\node[step, below=of ep, fill=orange!12] (rej){Amplitude reject \,---\, drop $>$\,150\,$\mu$V epochs};
\node[step, below=of rej]  (z) {Per-recording z-score \,---\, session invariance};

\foreach \a/\b in {n/bp, bp/car, car/ep, ep/rej, rej/z}
    \draw[arow] (\a) -- (\b);

\end{tikzpicture}
\caption{Stage 2: the signal processing sub-pipeline applied to
every recording. Each operation is applied identically during
training and inference. The output is a set of clean, normalised,
fixed-length epochs ready for the activity gate and feature
extraction.}
\label{fig:mo_preprocessing}
\end{figure}

The signal processing stage is deterministic and parameter-driven:
the same filters, the same window length, and the same normalisation
are applied to every epoch the system ever processes. The detailed
filter mathematics, zero-phase filtering rationale, epoch overlap
strategy, and the consequences for evaluation are presented in
Chapter~7.

%% ----------------------------------------------------------
\section{Stage 3: Activity Gating}
\label{sec:mo_gating}
%% ----------------------------------------------------------

Before features are extracted, the activity gate examines each
epoch and decides whether it contains a genuine gesture or merely a
rest window. This single decision is one of the most consequential
in the entire pipeline: it removes the mislabelled rest epochs that
would otherwise teach the classifier contradictory lessons, and it
does so using the same physiological channels that carry the
gestures themselves.

\begin{figure}[htbp]
\centering
\begin{tikzpicture}[
    node distance=1.0cm,
    box/.style={rectangle, rounded corners=2pt, draw=black!60, thick,
                minimum width=2.6cm, minimum height=0.9cm,
                align=center, font=\scriptsize},
    epoch/.style={box, fill=gray!10},
    keep/.style={box, fill=green!15},
    drop/.style={box, fill=red!12},
    dec/.style={diamond, aspect=2, draw=black!70, very thick,
                fill=blue!10, align=center, font=\scriptsize,
                inner sep=1pt},
    arow/.style={-{Stealth[length=2mm]}, thick},
]

\node[epoch] (ep) {Epoch\\\scriptsize 19 $\times$ 400};
\node[dec, right=1.4cm of ep] (d)
    {RMS of Fp1/Fp2\\or T3/T4\\$\geq$ 20\,$\mu$V?};
\node[keep, above right=0.4cm and 1.6cm of d] (k)
    {Keep\\\scriptsize active gesture};
\node[drop, below right=0.4cm and 1.6cm of d] (dr)
    {Discard\\\scriptsize rest window};

\draw[arow] (ep) -- (d);
\draw[arow] (d) -- node[above, font=\scriptsize] {yes} (k);
\draw[arow] (d) -- node[below, font=\scriptsize] {no}  (dr);

\end{tikzpicture}
\caption{Stage 3: the activity gate. Each epoch's RMS energy at the
gesture-relevant electrodes (Fp1/Fp2 for blinks, T3/T4 for jaw) is
compared against a 20\,$\mu$V threshold. Active epochs proceed to
feature extraction; rest epochs --- roughly half of all epochs ---
are discarded, removing the mislabelled-rest problem and yielding a
+3.1\% accuracy improvement.}
\label{fig:mo_gating}
\end{figure}

Across the dataset, the gate retains roughly half of all extracted
epochs, discarding the rest windows so that only genuine gesture
epochs are used for training and evaluation. The gate's design,
threshold calibration, and quantitative impact are detailed in
Chapter~7 alongside the rest of the signal processing pipeline.

%% ----------------------------------------------------------
\section{Stage 4: Feature Extraction}
\label{sec:mo_features}
%% ----------------------------------------------------------

Each surviving epoch --- a $19 \times 400$ matrix of cleaned samples
--- is converted into a channel-aware feature vector. Rather
than applying a uniform transform across all channels, the system
extracts different feature sets from different electrode groups,
each chosen for its physiological relevance to the artifacts being
classified.

\begin{figure}[htbp]
\centering
\begin{tikzpicture}[
    node distance=1.0cm,
    epoch/.style={rectangle, rounded corners=2pt, draw=black!70,
                  very thick, minimum width=2.4cm, minimum height=1.0cm,
                  align=center, font=\scriptsize, fill=gray!12},
    grp/.style={rectangle, rounded corners=2pt, draw=black!60, thick,
                minimum width=3.0cm, minimum height=1.1cm,
                align=center, font=\scriptsize},
    frontal/.style={grp, fill=blue!12},
    temporal/.style={grp, fill=orange!15},
    global/.style={grp, fill=green!12},
    vec/.style={rectangle, rounded corners=2pt, draw=black!70,
                very thick, minimum width=2.6cm, minimum height=1.0cm,
                align=center, font=\scriptsize, fill=purple!12},
    arow/.style={-{Stealth[length=2mm]}, thick},
]

\node[epoch] (ep) {Epoch\\\scriptsize 19 $\times$ 400};

\node[frontal, right=2.0cm of ep, yshift=1.4cm] (f)
    {Frontal (Fp1, Fp2)\\\scriptsize blink shape\\\scriptsize 38 features};
\node[temporal, right=2.0cm of ep] (t)
    {Temporal (T3, T4)\\\scriptsize jaw EMG + asymmetry\\\scriptsize 36 features};
\node[global, right=2.0cm of ep, yshift=-1.4cm] (g)
    {Global (19 ch)\\\scriptsize kurtosis\\\scriptsize 19 features};

\node[vec, right=2.2cm of t] (v)
    {Feature vector\\\scriptsize 93 dimensions};

\draw[arow] (ep.east) -- (f.west);
\draw[arow] (ep.east) -- (t.west);
\draw[arow] (ep.east) -- (g.west);
\draw[arow] (f.east) -- (v.west);
\draw[arow] (t.east) -- (v.west);
\draw[arow] (g.east) -- (v.west);

\end{tikzpicture}
\caption{Stage 4: channel-aware feature extraction. The epoch is
split into three physiologically motivated channel groups. The
frontal group captures blink morphology, the temporal group
captures jaw EMG and the left/right asymmetry that separates
bruxism directions, and the global group captures whole-scalp
statistics. The three subsets concatenate into one 93-dimensional
vector.}
\label{fig:mo_features}
\end{figure}

This channel-aware design directly encodes domain knowledge: blink
information lives at the front of the head, jaw information lives at
the temporal sites, and the difference between left and right
bruxism is captured by an explicit asymmetry index. The complete
feature inventory, all defining equations, and the feature
importance analysis are presented in Chapter~8.

%% ----------------------------------------------------------
\section{Stage 5: Classification}
\label{sec:mo_classification}
%% ----------------------------------------------------------

The feature vector is passed to the \textbf{hybrid classification
architecture} that assigns it to one of the five gesture classes.
Rather than a single monolithic classifier, the system uses a
two-stage router matched to the two signal families: hand-crafted
features first separate blink gestures from muscle gestures, and each
family is then sub-classified by the method best suited to its
physiology --- Dynamic Time Warping shape-matching for blinks, and
Riemannian covariance geometry for muscle gestures. During inference,
this router predicts a class (and a confidence score) for each
incoming epoch.

\begin{figure}[htbp]
\centering
\begin{tikzpicture}[
    node distance=1.0cm,
    box/.style={rectangle, rounded corners=2pt, draw=black!60, thick,
                minimum width=2.6cm, minimum height=0.9cm,
                align=center, font=\scriptsize},
    data/.style={box, fill=gray!10},
    proc/.style={box, fill=blue!10},
    dtw/.style={box, fill=green!12},
    riem/.style={box, fill=orange!15},
    out/.style={box, fill=purple!12},
    dec/.style={diamond, aspect=1.6, draw=black!70, thick,
                fill=blue!10, align=center, font=\scriptsize, inner sep=1pt},
    arow/.style={-{Stealth[length=2mm]}, thick},
]

\node[data] (v) {Feature\\vector};
\node[dec, right=1.2cm of v] (route) {blink\\or\\muscle?};
\node[dtw, above right=0.5cm and 1.3cm of route] (dtw) {DTW--KNN\\\scriptsize blink shape};
\node[riem, below right=0.5cm and 1.3cm of route] (riem) {Riemannian\\\scriptsize muscle covariance};
\node[out, below right=1.4cm and 1.3cm of dtw] (cls) {Class + conf.\\\scriptsize 1 of 5};

\draw[arow] (v) -- (route);
\draw[arow] (route) -- node[above, font=\scriptsize]{blink} (dtw);
\draw[arow] (route) -- node[below, font=\scriptsize]{muscle} (riem);
\draw[arow] (dtw) -- (cls);
\draw[arow] (riem) -- (cls);

\end{tikzpicture}
\caption{Stage 5: the hybrid classifier. Hand-crafted features first
route each epoch to the blink or muscle family; blink epochs are then
sub-classified by Dynamic Time Warping shape-matching, and muscle
epochs by Riemannian covariance geometry. Under the strict
cross-session evaluation protocol, the full architecture reports
91.3\% accuracy with a macro-F1 of 0.913.}
\label{fig:mo_classification}
\end{figure}

The design of this hybrid architecture, the failed alternatives that
motivated it, the \acrshort{smote} balancing protocol, and ---
critically --- the data leakage audit and cross-session evaluation
that distinguish honest from inflated accuracy figures are all
covered in detail in Chapter~9.

%% ----------------------------------------------------------
\section{Stage 6: Real-Time Inference and Actuation}
\label{sec:mo_inference}
%% ----------------------------------------------------------

In live operation, a continuous \acrshort{eeg} stream is processed
through a sliding window. Each window is preprocessed, gated,
feature-extracted, and classified exactly as in training. To prevent
single misclassifications from producing erratic hand movements, the
raw predictions pass through two safety layers: a confidence
threshold that rejects uncertain predictions, and a Finite State
Machine that requires several consecutive identical predictions
before a command is issued.

\begin{figure}[htbp]
\centering
\begin{tikzpicture}[
    node distance=0.75cm,
    box/.style={rectangle, rounded corners=2pt, draw=black!60, thick,
                minimum width=6.2cm, minimum height=0.7cm,
                align=center, font=\scriptsize, fill=green!10},
    safe/.style={box, fill=orange!12},
    hw/.style={box, fill=red!12},
    arow/.style={-{Stealth[length=2mm]}, thick},
]

\node[box] (win) {Sliding window \,---\, 2.0\,s, 0.2\,s step};
\node[box, below=of win] (proc) {Preprocess + gate + features \,(shared)};
\node[box, below=of proc] (pred) {Hybrid classifier prediction};
\node[safe, below=of pred] (conf) {Confidence threshold \,$\geq$ 0.60};
\node[safe, below=of conf] (fsm) {FSM \,---\, 3 consecutive agree};
\node[hw, below=of fsm] (ser) {UART label \,$\rightarrow$ ESP32};
\node[hw, below=of ser] (act) {Servo actuation \,---\, hand moves};

\foreach \a/\b in {win/proc, proc/pred, pred/conf, conf/fsm, fsm/ser, ser/act}
    \draw[arow] (\a) -- (\b);

\end{tikzpicture}
\caption{Stage 6: real-time inference and actuation. The shared
preprocessing and feature stages (green) mirror the training path
exactly. Two safety layers (orange) --- a confidence threshold and a
three-window FSM agreement requirement --- filter transient
misclassifications before a single command label is sent over UART to
the ESP32 controller, which drives the servo motors.}
\label{fig:mo_inference}
\end{figure}

Only when both safety conditions are satisfied is a single command
label transmitted over the UART link to the ESP32 controller, which
resolves it to a hand action and drives the servo motors. The sliding
window mechanics, threshold calibration, FSM design, UART command
protocol, and end-to-end latency analysis are presented in
Chapter~12.

%% ----------------------------------------------------------
\section{Chapter Summary}
\label{sec:mo_summary}
%% ----------------------------------------------------------

This chapter provided a high-level overview of the complete
methodology, tracing the path from voluntary artifact gesture to
prosthetic hand action through six pipeline stages. The system
converts a continuous 19-channel \acrshort{eeg} recording into a
single command label by progressively cleaning, windowing, gating,
and summarising the signal, then classifying the result through the
hybrid architecture and filtering the predictions through safety
layers before actuation. The defining architectural principle is that
the preprocessing, gating, and feature extraction stages are shared
identically between offline training and online inference,
guaranteeing that the model sees consistent data in both modes.

Table~\ref{tab:mo_chapter_map} maps each pipeline stage to the
chapter in which it is described in full detail.

\begin{table}[htbp]
\centering
\caption{Mapping of pipeline stages to their detailed chapters.}
\label{tab:mo_chapter_map}
\renewcommand{\arraystretch}{1.3}
\begin{tabular}{l l l}
\toprule
\textbf{Stage} & \textbf{Description} & \textbf{Detailed in} \\
\midrule
Stage 1    & EEG acquisition (hardware)   & Chapter~4 \\
Stage 1--2 & Data collection protocol     & Chapter~6 \\
Stage 2    & Signal processing            & Chapter~7 \\
Stage 3    & Activity gating              & Chapter~7 \\
Stage 4    & Feature extraction           & Chapter~8 \\
Stage 5    & Classification (hybrid)      & Chapter~9 \\
Stage 5    & Results \& analysis           & Chapter~10 \\
Stage 6    & Real-time inference          & Chapter~12 \\
\bottomrule
\end{tabular}
\end{table}

With this mental map in place, the following chapters examine each
stage in depth, beginning with the data collection protocol in
Chapter~6.